\documentclass[10pt,a4paper]{article}
\usepackage{amsmath,amssymb,amsthm,graphicx}
\usepackage{amsmath,amssymb,amsthm,amsfonts}
\usepackage{latexsym,bbm,graphicx,float,mathtools}
\usepackage{subcaption}

\usepackage{verbatim}
\usepackage{enumitem}
\usepackage{hyperref}  
\hypersetup{
    colorlinks,
    citecolor=black,
    filecolor=black,
    linkcolor=black,
    urlcolor=black,
    linktoc=all
}

\usepackage[noend]{algpseudocode}
\usepackage{algorithm}

\usepackage{thmtools}
\usepackage{thm-restate}

%primed-thm
\newtheorem{manualtheoreminner}{Theorem}
\newenvironment{manualtheorem}[1]{%
  \renewcommand\themanualtheoreminner{#1}%
  \manualtheoreminner
}{\endmanualtheoreminner}

%footnote twice preamble
%https://tex.stackexchange.com/questions/35043/reference-different-places-to-the-same-footnote
\makeatletter
\newcommand\footnoteref[1]{\protected@xdef\@thefnmark{\ref{#1}}\@footnotemark}
\makeatother



 % environments
\newtheorem{problem}{Problem}
\newtheorem{theorem}{Theorem}[section]
\newtheorem{lem}[theorem]{Lemma}

\newtheorem{proposition}[theorem]{Proposition}
\newtheorem{fact}[theorem]{Fact}
\newtheorem{cor}[theorem]{Corollary}
\newtheorem{conjecture}[theorem]{Conjecture}
\newtheorem{question}{Question}
\newtheorem{goal}{Goal}

% non-italicized environments
\theoremstyle{definition}%\newtheorem{definition}[theorem]{Definition}
%\newtheorem{remark}[theorem]{Remark}
\newtheorem{remark}[theorem]{Remark}
\newtheorem{claim}[theorem]{Claim}
\newtheorem{definition}[theorem]{Definition}
\newtheorem{obs}[theorem]{Observation}
\newtheorem{notation}[theorem]{Notation}
\newtheorem{example}[theorem]{Example}
\newtheorem{examples}[theorem]{Examples}
\theoremstyle{plain}

% probability
\newcommand{\std}{\operatorname{{\bf std}}}
\newcommand{\Var}{\operatorname{{\bf Var}}}
\newcommand{\Cov}{\operatorname{{\bf Cov}}}
\newcommand{\Varx}{\mathop{{\bf Var}\/}}
%\newcommand{\E}{\operatorname{{\mathbb E}}}
\DeclareMathOperator*{\E}{\mathbb{E}}
%\renewcommand{\Pr}{\operatorname{{\mathbb P}}}
\let\Pr\relax
\DeclareMathOperator*{\Pr}{\mathbb{P}}
\newcommand{\Bin}{\operatorname{Bin}}
\newcommand{\Ber}{\operatorname{Ber}}

% mathrm terms
\newcommand{\sbp}{\sum\mathrm{smABP}}
\newcommand{\poly}{\mathrm{poly}}
\newcommand{\imm}[2]{\mathrm{IMM}_{#1,#2}}
\newcommand{\rk}{\mathrm{relrk}}
\newcommand{\polylog}{\mathrm{polylog}}
\newcommand{\size}{\mathrm{size}}
\newcommand{\avg}{\mathop{\mathrm{avg}}}
\newcommand{\odd}{\mathrm{odd}}
\newcommand{\even}{\mathrm{even}}
\newcommand{\sign}{\mathrm{sign}}
\newcommand{\dist}{\mathrm{dist}}
\newcommand{\dtv}{d_\mathrm{TV}}
\newcommand{\supp}{\mathrm{supp}}
\newcommand{\Vol}{\mathrm{Vol}}
\newcommand{\opt}{\mathrm{opt}}
\newcommand{\rank}{\mathrm{rank}}
\newcommand{\Tr}{\mathrm{Tr}}
\newcommand{\Alg}{\mathrm{Alg}}
\newcommand{\sen}{\mathrm{sen}}
\newcommand{\dnf}{\mathrm{dnf}}
\newcommand{\dt}{\mathrm{dt}}
\newcommand{\type}{\mathrm{type}}
\newcommand{\cent}{\mathrm{center}}
\newcommand{\dyck}{\mathrm{Dyck}}
\newcommand{\out}{\mathrm{out}}
\newcommand{\children}{\mathrm{children}}
\newcommand{\Cop}{\mathsf{Copies}}

%mathsf complexity classes
\newcommand{\VP}{\mathsf{VP}}
\newcommand{\VBP}{\mathsf{VBP}}
\newcommand{\VNP}{\mathsf{VNP}}
\newcommand{\VF}{\mathsf{VF}}
%\newcommand{\P}{\mathsf{P}}
\newcommand{\NP}{\mathsf{NP}}




% mathbf terms
\newcommand{\Size}{\mathbf{Size}}
\newcommand{\Depth}{\mathbf{Depth}}

% number systems
\newcommand{\R}{\mathbb R}
\newcommand{\C}{\mathbb C}
\newcommand{\N}{\mathbb N}
\newcommand{\Z}{\mathbb Z}
\newcommand{\F}{\mathbb F}
\newcommand{\K}{\mathbb K}
\newcommand{\Q}{\mathbb Q}

% short forms
\newcommand{\eps}{\varepsilon}
\renewcommand{\epsilon}{\eps} 
\newcommand{\la}{\langle}
\newcommand{\ra}{\rangle}
\newcommand{\wt}[1]{\widetilde{#1}}
\newcommand{\wh}[1]{\widehat{#1}}
\newcommand{\sse}{\subseteq}
\newcommand{\bits}{\{-1,1\}}
\newcommand{\zo}{\{0,1\}}
\newcommand{\bn}{\bits^n}
\newcommand{\isafunc}{:\bn\to\bits}
\newcommand{\isazofunc}{:\zo^n\to\zo}
\newcommand{\sumi}{\sum_{i=1}^n}

\newcommand{\Fsm}{\mathbb{F}_{\mathrm{sm}}}
\newcommand{\Ksm}{\mathbb{K}_{\mathrm{sm}}}
\newcommand{\lfrac}{\textstyle \frac}



% Overline
\newcommand{\ol}[1]{\overline{#1}}

% BOLD
\newcommand{\ind}{{\bf 1}}

\newcommand{\bA}{\mathbf{A}}
\newcommand{\bB}{\mathbf{B}}
\newcommand{\bC}{\mathbf{C}}
\newcommand{\bD}{\mathbf{D}}
\newcommand{\bE}{\mathbf{E}}
\newcommand{\bF}{\mathbf{F}}
\newcommand{\bG}{\mathbf{G}}
\newcommand{\bH}{\mathbf{H}}
\newcommand{\bI}{\mathbf{I}}
\newcommand{\bJ}{\mathbf{J}}
\newcommand{\bK}{\mathbf{K}}
\newcommand{\bL}{\mathbf{L}}
\newcommand{\bM}{\mathbf{M}}
\newcommand{\bN}{\mathbf{N}}
\newcommand{\bO}{\mathbf{O}}
\newcommand{\bP}{\mathbf{P}}
\newcommand{\bQ}{\mathbf{Q}}
\newcommand{\bR}{\mathbf{R}}
\newcommand{\bS}{\mathbf{S}}
\newcommand{\bT}{\mathbf{T}}
\newcommand{\bU}{\mathbf{U}}
\newcommand{\bV}{\mathbf{V}}
\newcommand{\bW}{\mathbf{W}}
\newcommand{\bX}{\mathbf{X}}
\newcommand{\bY}{\mathbf{Y}}
\newcommand{\bZ}{\mathbf{Z}}


\newcommand{\boldf}{\boldsymbol{f}} % \bf taken
\newcommand{\ba}{\boldsymbol{a}}
\newcommand{\bb}{\boldsymbol{b}}
\newcommand{\be}{\boldsymbol{e}}
\newcommand{\bg}{\boldsymbol{g}}
\newcommand{\bh}{\boldsymbol{h}}
\newcommand{\bj}{\boldsymbol{j}}
\newcommand{\bk}{\boldsymbol{k}}
\newcommand{\bm}{\boldsymbol{m}}
\newcommand{\bp}{\boldsymbol{p}}
\newcommand{\boldr}{\boldsymbol{r}} % \br taken
\newcommand{\bs}{\boldsymbol{s}}
\newcommand{\bt}{\boldsymbol{t}}
\newcommand{\bu}{\boldsymbol{u}}
\newcommand{\bv}{\boldsymbol{v}}
\newcommand{\bw}{\boldsymbol{w}}
\newcommand{\bx}{\boldsymbol{x}}
\newcommand{\by}{\boldsymbol{y}}
\newcommand{\bz}{\boldsymbol{z}}

\newcommand{\boldR}{\boldsymbol{R}}


\newcommand{\bsigma}{\boldsymbol{\sigma}}
\newcommand{\bpi}{\boldsymbol{\pi}}
\newcommand{\bell}{\boldsymbol{\ell}}
\newcommand{\bnu}{\boldsymbol{\nu}}
\newcommand{\bxi}{\boldsymbol{\xi}}
\newcommand{\btau}{\boldsymbol{\tau}} 
\newcommand{\brho}{{\boldsymbol{\rho}} }
\newcommand{\bgamma}{{\boldsymbol{\gamma}} }
\newcommand{\bphi}{{\boldsymbol{\varphi}} }


\newcommand{\bvsigma}{{\boldsymbol{\varsigma}} }
\newcommand{\bvrho}{{\boldsymbol{\varrho}}}

\newcommand{\bGamma}{{\mathbf{\Gamma}} }


% Calligraphy

\newcommand{\calA}{\mathcal{A}}
\newcommand{\calB}{\mathcal{B}}
\newcommand{\calC}{\mathcal{C}}
\newcommand{\calD}{\mathcal{D}}
\newcommand{\calE}{\mathcal{E}}
\newcommand{\calF}{\mathcal{F}}
\newcommand{\calG}{\mathcal{G}}
\newcommand{\calH}{\mathcal{H}}
\newcommand{\calI}{\mathcal{I}}
\newcommand{\calJ}{\mathcal{J}}
\newcommand{\calK}{\mathcal{K}}
\newcommand{\calL}{\mathcal{L}}
\newcommand{\calM}{\mathcal{M}}
\newcommand{\calN}{\mathcal{N}}
\newcommand{\calO}{\mathcal{O}}
\newcommand{\calP}{\mathcal{P}}
\newcommand{\calQ}{\mathcal{Q}}
\newcommand{\calR}{\mathcal{R}}
\newcommand{\calS}{\mathcal{S}}
\newcommand{\calT}{\mathcal{T}}
\newcommand{\calU}{\mathcal{U}}
\newcommand{\calV}{\mathcal{V}}
\newcommand{\calW}{\mathcal{W}}
\newcommand{\calX}{\mathcal{X}}
\newcommand{\calY}{\mathcal{Y}}
\newcommand{\calZ}{\mathcal{Z}}

\newenvironment{solution}
  {\begin{proof}[\textbf{Solution}]}
  {\end{proof}}
 
\usepackage{cleveref}
\crefname{originalsection}{original section}{original sections}
\Crefname{originalsection}{Original section}{Original sections}
\makeatletter
\newlabel{subsec:our-results}{{1.3}{source}{}{sourcepointer-1}{}}
\newlabel{subsec:our-results@cref}{{[originalsection][0][]1.3}{[1][0][]source}{}{}{}}
\makeatother
\title{Low-degree superpolynomial and high-degree exponential lower bounds for sums of ordered set-multilinear ABPs}
\author{Prerona Chatterjee, Deepanshu Kush, Shubhangi Saraf, Amir Shpilka}
\date{doi:10.4230/LIPIcs.CCC.2024.20; CC BY-NC-SA4.0}
\begin{document}
\maketitle
\noindent\textit{Editorial note.} Original source: Prerona Chatterjee, Deepanshu Kush, Shubhangi Saraf and Amir Shpilka; arXiv:2312.15874v3 (2024-02-19), https://arxiv.org/abs/2312.15874v3. License: Creative Commons Attribution-NonCommercial-ShareAlike4.0, https://creativecommons.org/licenses/by-nc-sa/4.0/.\par\medskip
\noindent\textit{Editorial note.} Non-commercial research collection. Changes: complete author theorem/proof/context excerpts and original macros/bibliography in standalone layout; unrelated sections and frontmatter omitted. No mathematical corrections or additional proofs. Adapted presentation retains CC BY-NC-SA4.0 including NC/SA restrictions.\par\medskip
\noindent\textit{Allowed standard prerequisites (editorial).}\begin{itemize}
\item Relative-rank properties and inner-product gadget. Author Claim clm:rk-props and Observation obs:ip-full-rank; exact statements retained.
\item Explicit VP hard polynomial with full relative rank. Author Lemma lem:Fn-VP-props explicitly cites KushS23; source construction and full quoted guarantee retained.
\item Standard multiplicative Chernoff bound. Author Theorem thm:Chernoff explicitly cites Mitzenmacher-Upfal; exact source statement retained.
\end{itemize}\medskip
\ifdefined\hypertarget\hypertarget{sourcepointer-1}{}\fi
\noindent\textit{Source reference.} Original section 1.3: number-theory/sources/ccc-2024-20/native/main.tex, source lines 150--150. Optional introductory pointer to results; counted full target statements and proofs are included.\par\medskip
An \emph{algebraic circuit} over a field $\F$ for a multivariate polynomial $P(x_1,\ldots,x_N)$ is a directed acyclic graph (DAG)
whose internal vertices (called gates) are labeled as either $+$ (sum) or $\times$ (product), and leaves
(vertices of in-degree zero) are labeled by the variables $x_i$ or constants from $\F$. 
A special output gate (the root of the DAG) represents the polynomial
$P$. 
If the DAG happens to be a tree, such a resulting circuit is called an \emph{algebraic formula}. 
The size of a circuit or formula is the number of nodes in the DAG. We also consider the product-depth of the circuit, which is the maximum number of product gates on a root-to-leaf path. 
The class $\VP$ (respectively, $\VF$) is then defined to be the collection of all polynomials having at most polynomially large degree which can be computed by polynomial-sized circuits (respectively, formulas).

An \emph{algebraic branching program} (ABP) is a layered DAG with two special nodes in
it: a start-node and an end-node. All edges of the ABP go from layer $\ell -1$ to layer $\ell$ for some $\ell$ (say start-node is the unique node in layer $0$ and end-node is the unique node in the last layer) and are labeled by a linear polynomial. Every directed path $\gamma$ from start-node to end-node computes the monomial $P_\gamma$,
which is the product of all labels on the path $\gamma$. The ABP computes the polynomial $P = \sum_{\gamma} P_\gamma$,
where the sum is over all paths $\gamma$ from start-node to end-node. Its size is simply the number of nodes in the DAG, its \emph{depth} is the length of the longest path from the start-node to the end-node, and \emph{width} is the maximum number of nodes in any layer. The class $\VBP$ is then defined to be the collection of all polynomials (with polynomially-bounded degree) which can be computed by polynomial-sized branching programs. ABPs are known to be of intermediate complexity between formulas and circuits; in other words, we know the inclusions $\VF\subseteq\VBP \subseteq \VP\sse \VNP$.

A polynomial is said to be \emph{homogeneous}
if each monomial has the same total degree and \emph{multilinear} if every variable occurs
at most once in any monomial. Now, suppose that the underlying variable set is partitioned into $d$ sets
$X_1,\ldots, X_d$. Then the polynomial is said to be \emph{set-multilinear} with respect to this variable partition if each
monomial in $P$ has \emph{exactly} one variable from each set. Note that a set-multilinear polynomial is both multilinear and homogeneous, and has degree precisely $d$ if it is set-multilinear over $d$ sets. Next, we define different models
of computation corresponding to these variants of polynomials classes. An algebraic formula/branching program/circuit is set-multilinear with respect to a variable partition
$(X_1,\ldots, X_d)$ if each internal node in the formula/branching program/circuit computes a set-multilinear polynomial.\footnote{Of course, a non-root node need not be set-multilinear with respect to the \emph{entire} variable partition. Nevertheless, here we demand that it must be set-multilinear with respect to some \emph{subset} of the collection $\{X_1,\ldots,X_d\}$.} Multilinear and homogeneous formulas/branching programs/circuits are defined analogously.  

\begin{definition}[Ordered smABP]\label{def:smABP}
Given a variable partition $(X_1,\ldots, X_d)$, we say that a set-multilinear branching program of depth $d$ is said to be \emph{ordered} with respect to an 
ordering (or permutation) $\sigma\in S_d$ if for each $\ell\in[d]$, all edges of the ABP from layer $\ell-1$ to layer $\ell$ are labeled using a linear form over the variables in $X_{\sigma(\ell)}$. It is simply said to be \emph{ordered} if there \emph{exists} an ordering $\sigma$ such that it is ordered with respect to $\sigma$.
\end{definition}

At this point, it is essential to take note of the terminology in this context: in this paper, a general (or ``unordered'') set-multilinear branching program refers to an ABP for which each internal node computes a polynomial that is set-multilinear with respect to \emph{some} subset of the global partition, whereas an \emph{ordered} set-multilinear branching program is more specialized and has the property that any two nodes in the same layer compute polynomials that are set-multilinear with respect to the \emph{same} partition.

\begin{definition}[$\sbp$]\label{def:sum}
Given a polynomial $P(X)$ that is set-multilinear with respect to the variable partition $X = (X_1,\ldots, X_d)$, we say that $\sum_{i=1}^t A_i$ is a $\sbp$ computing $P$ if indeed $\sum_{i=1}^t A_i(X) = P(X)$, and
each $A_i$ is an ordered set-multilinear branching program i.e., each $A_i$ is ordered with respect to some $\sigma_i\in S_d$. We call $t$ (i.e., the number of summands in a $\sbp$) its \emph{support size} and define its \emph{max-width} and \emph{total-width} to be the maximum over the width of each $A_i$ and the sum of the width of each $A_i$, respectively.    
\end{definition}

\begin{theorem}[``Low''-Degree $\sbp$ Lower Bounds]\label{thm:small-deg-VP}
Let $d\leq n$ be growing parameters satisfying $d = \omega(\log n)$. There is a $\Theta(dn)$-variate degree $d$ set-multilinear polynomial $F_{n,d}$ in $\VP$ such that $F_{n,d}$  cannot be computed by a $\sbp$ of total-width $\poly(n)$.
\end{theorem}

\begin{theorem}[Exponential Lower Bounds for $\sbp$]\label{thm:large-deg-VP}
There is a set-multilinear polynomial $F_{n,n}$ in $\VP$, in  $\Theta(n^2)$ variables and of degree $\Theta(n)$, such that any $\sbp$ computing $F_{n,n}$ requires total-width ${\exp(\Omega(n^{1/3}))}$.
\end{theorem}

\section{Preliminaries}\label{sec:prel}
\subsection{Relative Rank and its Properties}

We first describe the notation that we need to define the measures that we use to prove our results described in Section \ref{subsec:our-results}. Instead of directly working with the rank of the partial derivative matrix, we work with the following normalized form.

\begin{definition}\label{def:rk-rest}
Let $w = (w_1, w_2,\ldots, w_d)$ be a tuple (or word) of non-zero real numbers. For a subset $S \subseteq [t]$, we shall refer to the sum $\sum_{i\in S} w_i$ by $w_S$,  and by $w|_S$, we will refer to the tuple obtained by considering only the
elements of $w$ that are indexed by $S$. 
Given a word $w=(w_1,\ldots,w_d)$, we denote by $\ol{X}(w)$ a tuple of $d$ sets of variables $(X(w_1),\ldots,X(w_d))$ where
$|X(w_i)| = 2^{|w_i|}$.\footnote{In particular, $2^{|w_i|}\in\mathbb{N}$.}
We denote by $\Fsm[\calT]$ the set of set-multilinear polynomials over the tuple
of sets of variables $\calT$.
\end{definition}


\begin{definition}[Relative Rank Measure of \cite{LST1}]\label{def:relrk}
Let $\ol{X}=(X_1,\ldots,X_d)$ be a tuple of sets of variables such that $|X_i|=n_i$ and let $f\in \Fsm[\ol{X}]$.
 Let $w = (w_1, w_2,\ldots, w_d)$ be a tuple (or word) of non-zero real numbers such that $2^{|w_i|}  = n_i $ for all $i \in [d]$. Corresponding to a word $w$, define $\calP_w \coloneqq \{i\ |\ w_i > 0\}$ and $\calN_w \coloneqq \{i\ |\ w_i < 0\}$. Let $\calM^{\calP}_{w}$ be the
set of all set-multilinear monomials over the subset of the variable sets $X_1, X_2,\ldots, X_d$ precisely
indexed by $\calP_w$, and similarly let $\calM^{\calN}_{w}$ be the set of all set-multilinear monomials over these variable
sets indexed by $\calN_w$.

Define the ‘partial derivative matrix’ matrix $\calM_w(f)$ whose rows are indexed by the elements
of $\calM^{\calP}_w$ and columns indexed by the elements of $\calM^{\calN}_{w}$ as follows: the entry of this matrix corresponding to a
row $m_1$ and a column $m_2$ is the coefficient of the monomial $m_1\cdot m_2$ in $f$. We define
\[
\rk_w(f) \coloneqq \frac{\mathrm{rank}(\calM_w(f))}{\sqrt{|\calM^{\calP}_w|\cdot |\calM^{\calN}_w|}} = \frac{\mathrm{rank}(\calM_w(f))}{2^{\frac{1}{2}\sum_{i\in[d]}|w_i|}}.
\]
\end{definition}


The following is a simple result that establishes various useful properties of the relative rank measure.
\begin{claim}[\cite{LST1}]\label{clm:rk-props}
\begin{enumerate}
    \item(Imbalance) Say $f \in \Fsm[\ol{X}(w)]$. Then, $\rk_w(f)\leq 2^{-|w_{[d]}|/2}$.
    \item(Sub-additivity) If $f,g \in \Fsm[\ol{X}(w)]$, then $\rk_w(f+g)\leq \rk_w(f)+\rk_w(g)$.
    \item(Multiplicativity) Say $f = f_1 f_2\cdots f_t $ and assume that for each $i\in [t]$, $f_i\in \Fsm[\ol{X}(w|_{S_i})]$, where $(S_1, \ldots, S_t)$ is a partition of $[d]$. Then
    \[
    \rk_w(f) = \prod_{i\in[t]} \rk_{w|_{S_i}}(f_i).
    \]
\end{enumerate}
\end{claim}

We will repeatedly make use of the following.

\begin{theorem}[Chernoff Bound, as stated in \cite{mitzenmacher_upfal_2005}]\label{thm:Chernoff}
Suppose $X_1, \ldots, X_n$ are independent random variables taking values in $\zo$. Let $X$ denote their sum and let $\mu = \E[X]$ denote the expected value of the sum. Then for any $\delta > 0$,
\begin{eqnarray*} \Pr[X\ge (1+\delta)\mu] &\le& \exp\left(-\frac{\delta^2\mu}{2+\delta}\right), \text{ if }0\leq \delta, \\
  \Pr[X\le (1-\delta)\mu] &\le& \exp\left(-\frac{\delta^2\mu}{2}\right), \text{ if }0< \delta <1, \\ 
  \Pr[|X-\mu|\geq \delta \mu]&\leq& 2\cdot \exp\left(-\frac{\delta^2\mu}{3}\right),\text{ if }0< \delta <1.\\
\end{eqnarray*}
\end{theorem}

\subsection{Inner Product Gadget}
The following observation is used crucially to construct the hard polynomials in $\VP$ as well as $\VBP$. 
\begin{obs}[\cite{KushS23}]\label{obs:ip-full-rank}
    Let $n = 2^k$ and $X_1 = \{x_{1,1},\ldots,x_{1,n}\}$ and $X_2 = \{x_{2,1},\ldots,x_{2,n}\}$ be two disjoint sets of variables. Then, for any \emph{symmetric} word $w\in\{k,-k\}^2$ (i.e., where $w_1 + w_2 = 0$) and for the inner product `gadget' $ f = X_1\cdot X_2 = \sumi x_{1,i}x_{2,i}$, $\rk_w(f) = 1$ i.e., $\calM_w(f)$ is full-rank.
\end{obs}


\subsection{A Hard Set-multilinear Polynomial in $\VP$}\label{sec:VP}
Let $d$ be an even integer and let $X = (X_1,\ldots, X_{d})$ be a collection of sets of variables where each $|X_i| = n$, and similarly, let $Y = (Y_1,\ldots, Y_{d})$ be a distinct collection of sets of variables where each $|Y_i| = n$. We shall refer to the $Y$-variables as the \emph{auxiliary} variables. For $i$ and $j \in \{1,\ldots,d\}$, let $X_i \cdot X_j$ denote the inner-product quadratic form $\sum_{k = 1}^n x_{ik}x_{jk}$. Here, we shall assume that $X_i = \{x_{i,1},\ldots, x_{i,n}\}$ and $Y_i = \{y_{i,1},\ldots, y_{i,n}\}$.

For two integers $i \in \N$ and $j \in \N$, we denote $[i, j] = \{k \in \N : i \leq k \text{ and } k \leq j\}$ and call such a set an \emph{interval}. For every interval $[i, j] \subseteq [d]$, we define a polynomial $f_{i,j}(X,Y)\in \Fsm[X_i,\ldots,X_j,Y_i,\ldots, Y_j]$ as follows:

\[
f_{i,j} = \begin{cases}
y_{i,j}y_{j,i}(X_i\cdot X_j) &\text{if $j = i+1$}\\
0 &\text{if $j-i$ is even}\\
y_{i,j}y_{j,i}(X_i\cdot X_j)\cdot f_{i+1,j-1} + \sum_{r = i+1}^{j-1} f_{i,r}f_{r+1,j} &\text{otherwise}\\
\end{cases}
\]

These $f_{i,j}$ in present form were defined in \cite{KushS23}, but were in turn inspired from an earlier work of Raz and Yehudayof (\cite{RazY08}) in the multilinear context. \cite{KushS23} shows that they have the following full-rank property that will be instrumental for us.
\begin{lem}[\cite{KushS23}]\label{lem:Fn-VP-props}
    Let $n = 2^k$ and $d\leq n$ be an even integer. Over any field $\F$ of characteristic zero, the polynomial $F_{n,d} = f_{1,d}\in\Fsm[X,Y]$ as defined above satisfies the following: For any $w\in\{-k,k\}^d$ with $w_{[d]} = 0$, $\calM_w(F_{n,d})$ is full-rank when viewed as a matrix over the field $\F(Y)$, the field of rational functions over the $Y$ variables.
\end{lem}


\section{Separation between $\VP$ and $\sbp$}

In the theorem below, $F_{n,d}$ refers to the polynomial defined in Section \ref{sec:VP}. 
In this section, we prove \autoref{thm:large-deg-VP} and \autoref{thm:small-deg-VP} by first proving the following statement. 

\begin{lem}\label{lem:VP}
    Given large enough integers $d\leq n$, any $\sbp$ of max-width $s$ and support size $t$ computing $F_{n,d}(X_1,\ldots,X_d,Y_1\ldots,Y_d)$ must satisfy at least one of the following:
    \begin{itemize}
        \item either $t>e^{d/96}$,
        \item or $t\cdot s^q \geq n^{\frac{\sqrt{dq}}{20}}$, for any integer $q$ in the range $ [12\ln(2t\sqrt{d}), d/4]$.
    \end{itemize}  
\end{lem}

\begin{proof}
    Suppose that $t\leq e^{d/96}$ (so that the range in the theorem statement is indeed well-defined). 
    
    First, we observe that for any $\sbp$ computing $F_{n,d}(X,Y)$, we can view each summand as an ordered set-multilinear branching program with respect to only the $X$ variables. 
    In other words, by appropriately collapsing the layers labelled using the $Y$ variables, each summand is a set-multilinear branching program over the field $\F(Y)$ \emph{ordered} with respect to $(X_{\sigma(1)},\ldots, X_{\sigma(d)})$ for some permutation $\sigma\in S_d$. 
    It is easy to see that this collapsing process does not increase the width or the size of any summand in the branching program. 
    The edge labels do get altered however: the coefficients of the $X$ variables in any edge label can now be polynomials in the $Y$ variables (and therefore, field constants in $\F(Y)$).

    Let $A$ be such a set-multilinear branching program of width $s$ and depth $d$ that is \emph{ordered} with respect to $(X_1,\ldots, X_d)$\footnote{In general, $A$ may be ordered with respect to an arbitrary permutation $\sigma$, but the assumption that $\sigma$ is the identity permutation in the discussion that follows is without loss of generality.}. 
    Recall that this means that for each $\ell\in[d]$, all edges of the ABP from layer $\ell-1$ to layer $\ell$ are labeled using a linear form {in} $\mathit{X_\ell}$. 
Given a node $u$ in layer $i$ and a node $v$ in layer $j>i$ of $A$, define $g_{u,v}$ to be the polynomial computed by the ABP restricted to the layers $i+1,\ldots,j-1$ with the source and the sink defined by $u$ and $v$ respectively. Consider an integer $q$ in the range specified in the theorem statement and let $r$ be the largest  integer $r$ such that the product $qr< d$ i.e., we must have $d-qr\leq q$. Consider the following decomposition of $A$:
\[
A(X_1,\ldots,X_d) = \sum_{u_1,\ldots, u_{q}} \prod_{i = 1}^{q+1} g_{u_{i-1},u_i},
\]
where $u_0$ and $u_{q+1}$ are defined to be the source and the sink of $A$ respectively, and for $1\leq i \leq q$, $u_i$ varies over all choices of nodes in layer $r\cdot i$. Note that hence, this expression contains at most $s^q$ terms. Also, note that each $g_{u_{i-1},u_i}$ is set-multilinear over the partition $(X_i)_{i\in S_i}$ where $S_i$ is the set $\{r(i-1) + 1, \ldots, ri\}$\footnote{If $A$ is instead ordered with respect to $\sigma$, then $S_i$ is taken to be the set $\{r(i-1) + 1, \ldots, ri\}$.} of length exactly $r$ for $i\in [q]$, and $S_{q+1}=\{rq+1,\ldots,d\}$ has length at most $q$. We now analyze the relative rank of each summand. 

Let  $w\in \{-k,k\}^d$ (where $k = \log n$) be an arbitrary word. By \autoref{clm:rk-props}, we see that 

\[
\rk_w\left(\prod_{i = 1}^{q+1} g_{u_{i-1},u_i}\right) = \prod_{i = 1}^{q+1} \rk_{w|_{S_i}}(g_{u_{i-1},u_i})
\leq \prod_{i = 1}^{q} \rk_{w|_{S_i}}(g_{u_{i-1},u_i})\leq \prod_{i = 1}^q 2^{-|w_{S_i}|/2}= 2^{-\frac12\left(\sum_{i=1}^q |w_{S_i}|\right)},
\]
from which we observe that the task of upper bounding this rank can be reduced to the task of lower bounding the sum $\sum_{i=1}^{q}|w_{S_i}|$, which is established below. 

Choose $w$ from $\{-k,k\}^d$ (where $k = \log n$) uniformly at random. For each $i\in[q]$, let $E_i$ denote the (bad) event that $|w_{S_i}|\leq  \sqrt{r}k/4$. Since $S_i$ is an interval of length $r$, by a standard estimation of binomial coefficients,  
%upon applying the Chernoff bound (\autoref{thm:Chernoff}), 
we obtain that $\Pr[E_i]\leq 1/4$. Then,  by the Chernoff bound (\autoref{thm:Chernoff}), the probability that at least half of the events $E_i$ occur is at most $e^{-\frac{q}{12}}$. Therefore, with probability at least $1 - e^{-\frac{q}{12}}$, 
\[
\sum_{i=1}^q |w_{S_i}| \geq {q\sqrt{r}k/8},
\]
and therefore, by the sub-additivity of $\rk_w(\cdot)$,
\[
\rk_w(A) \leq s^q 2^{-\frac{kq\sqrt{r}}{16}} 
= s^q n^{-\frac{{q\sqrt{r}}}{16}}.
\]

Now, let $\sum_{i=1}^t A_i$ be a $\sbp$ computing $F_{n,d}$ with max-width bounded by $s$, and such that each $A_i$ is ordered set-multilinear with respect to the variable partition $(X_{\sigma_i(1)},\ldots, X_{\sigma_i(d)})$ for some permutation $\sigma_i\in S_d$. %such that $\sum_{i=1}A_i$ computes a polynomial $P(X_1,\ldots, X_d)$ such that $\rk_w(P) = 1$ for every word $w\in\{-k,k\}^d$ with $w_{[d]} = 0$. %Such polynomials exist and we have explicit examples in VNP and VP due to \cite{KushS22} and \cite{KushS23} respectively.
By the union bound and the discussion above, it follows that with probability at least $1-t\cdot e^{-\frac{q}{12}}$, 
\[
\rk_w(F_{n,d}) = \rk_w\left(\sum_{i=1}^t A_i\right) \leq t\cdot s^q n^{-\frac{{q\sqrt{r}}}{16}}.
\]

But now, we can condition on the event that $w_{[d]} = 0$ (which occurs with probability $\Theta(\frac{1}{\sqrt{d}}$)) to establish the existence of a word $w\in \{-k,k\}^d$ with $w_{[d]} = 0$ such that $w$ satisfies $
\rk_w(P)\leq t\cdot s^q n^{-q\sqrt{r}}$. This is because of the given bound $q\geq 12\ln(2t\sqrt{d})$. Because $\rk_w(F_{n,d}) = 1$ for such a $w$ by \autoref{lem:Fn-VP-props}, we conclude that $t\cdot s^q \geq n^{\frac{{q\sqrt{r}}}{16}} \geq n^{\frac{\sqrt{(d-q)q}}{16}}\geq n^{\frac{\sqrt{dq}}{20}}$, where the last inequality follows from our choice of $q$.
%t\geq \left(\frac{n^{\frac{\alpha\sqrt{r}}{4}}}{s}\right)^q.
    
\end{proof}


\begin{proof}[Proof of \autoref{thm:large-deg-VP}]
We invoke Lemma \ref{lem:VP} with $d = n$. If $t \geq 2^{n^{1/3}}$, then we trivially have that the total-width is at least  ${\exp(\Omega(n^{1/3}))}$, so assume $t \leq 2^{n^{1/3}}$. We shall show that then, $s = \exp(\Omega(n^{1/3}))$, which will yield the desired result.

Set $q = \lceil 15n^{1/3}\rceil$. Then clearly, $q\leq n/4$. Moreover, as $t \leq 2^{n^{1/3}}$ by assumption, we verify that 
\[
q = \lceil 15n^{1/3} \rceil> 12 (n^{1/3} + \ln(2\sqrt{d})) \geq 12\ln(2t\sqrt{n}).
\]

Therefore, we can use \autoref{lem:VP} to obtain the inequality $t\cdot s^q \geq n^{\frac{\sqrt{nq}}{20}}$ (as $t\leq 2^{n^{1/3}}< e^{n/96}$ for large enough $n$). Plugging in $q = \lceil 15n^{1/3} \rceil$, we see that 
\[
s \geq \frac{n^{\frac{\sqrt{n/q}}{20}}}{t^{1/q}} \geq \frac{n^{\frac{n^{1/3}}{80}}}{2^{1/15}} = \exp(\Omega(n^{1/3})).
\]
\end{proof}

\begin{proof}[Proof of \autoref{thm:small-deg-VP}]
We consider cases as follows:

\paragraph{\underline{Case $t = \poly(n)$:}} Suppose there is a constant $c$ such that $t\leq n^c$. Set $q = 20 \ln (2n^c\sqrt{d})$ and note that $q = \Theta(\log n)$. We see that by \autoref{lem:VP},
\[
s \geq \frac{n^{\frac{\sqrt{d/q}}{20}}}{t^{1/q}} \geq n^{\frac{1}{20}\sqrt{\frac{d}{q}} - \frac{c}{q}}.
\]
Note that $c/q <1$ and decays to zero as $d$ becomes larger. Furthermore, as $d = \omega(\log n)$ by assumption and $q = \Theta(\log n)$, $d/q = \omega(1)$. We conclude that if $t$ is bounded by a polynomial in $n$, then $s$ must be super-polynomial in $n$.

\paragraph{\underline{Case $s = \poly(n)$:}} Suppose there is a fixed constant $c\geq 1$ such that $s\leq n^c$. Assume that $t<e^{d/96c^2}$, and set $q = d/1600c^2$. Then $q$ indeed lies in the range to apply \autoref{lem:VP}. We obtain the inequality
\[
t\geq \left(\frac{n^{\frac{\sqrt{d/q}}{20}}}{s}\right)^q \geq n^{\frac{\sqrt{dq}}{20} - cq} = n^{cq} = n^{\Omega(d)},
\]
which contradicts the assumption that $t<e^{d/96c^2}$. Hence, $t \geq e^{d/96c^2} = 2^{\Omega(d)}$ which is indeed super-polynomial in $n$ whenever $d = \omega(\log n)$.

Thus, in either case, it is shown that both $s$ and $t$ cannot be polynomially bounded. Hence, the total-width of any $\sbp$ computing $F_{n,d}$ cannot be polynomially bounded.
\end{proof}

\newcommand{\etalchar}[1]{$^{#1}$}
\begin{thebibliography}{DMPY12}

\bibitem[AFS{\etalchar{+}}18]{AndersonFSSV18}
Matthew Anderson, Michael~A. Forbes, Ramprasad Saptharishi, Amir Shpilka, and
  Ben~Lee Volk.
\newblock Identity testing and lower bounds for read-\emph{k} oblivious
  algebraic branching programs.
\newblock {\em {ACM} Trans. Comput. Theory}, 10(1):3:1--3:30, 2018.

\bibitem[AGKS15]{AgrawalGKS15}
Manindra Agrawal, Rohit Gurjar, Arpita Korwar, and Nitin Saxena.
\newblock Hitting-sets for {ROABP} and sum of set-multilinear circuits.
\newblock {\em {SIAM} J. Comput.}, 44(3):669--697, 2015.

\bibitem[AR16]{ArvindR16}
Vikraman Arvind and S.~Raja.
\newblock Some lower bound results for set-multilinear arithmetic computations.
\newblock {\em Chic. J. Theor. Comput. Sci.}, 2016, 2016.

\bibitem[BDS22]{mfcs/BhargavDS22}
C.~S. Bhargav, Sagnik Dutta, and Nitin Saxena.
\newblock Improved lower bound, and proof barrier, for constant depth algebraic
  circuits.
\newblock In Stefan Szeider, Robert Ganian, and Alexandra Silva, editors, {\em
  47th International Symposium on Mathematical Foundations of Computer Science,
  {MFCS} 2022, August 22-26, 2022, Vienna, Austria}, volume 241 of {\em
  LIPIcs}, pages 18:1--18:16. Schloss Dagstuhl - Leibniz-Zentrum f{\"{u}}r
  Informatik, 2022.

\bibitem[BDS24]{BDS23}
C.S. Bhargav, Prateek Dwivedi, and Nitin Saxena.
\newblock Lower bounds for the sum of small-size algebraic branching programs.
\newblock {\em To appear in the proceedings of the Annual Conference on Theory
  and Applications of Models of Computation}, 2024.
\newblock \url{https://www.cse.iitk.ac.in/users/nitin/papers/sumRO.pdf}.

\bibitem[BG22]{BhargavaG22}
Vishwas Bhargava and Sumanta Ghosh.
\newblock Improved hitting set for orbit of roabps.
\newblock {\em Comput. Complex.}, 31(2):15, 2022.

\bibitem[BS21]{BishtS21}
Pranav Bisht and Nitin Saxena.
\newblock Blackbox identity testing for sum of special roabps and its border
  class.
\newblock {\em Comput. Complex.}, 30(1):8, 2021.

\bibitem[B{\"{u}}r00]{burg}
Peter B{\"{u}}rgisser.
\newblock Cook's versus valiant's hypothesis.
\newblock {\em Theor. Comput. Sci.}, 235(1):71--88, 2000.

\bibitem[CKSV22]{ChatterjeeKSV22}
Prerona Chatterjee, Mrinal Kumar, Adrian She, and Ben~Lee Volk.
\newblock Quadratic lower bounds for algebraic branching programs and formulas.
\newblock {\em Comput. Complex.}, 31(2):8, 2022.

\bibitem[DMPY12]{DvirMPY12}
Zeev Dvir, Guillaume Malod, Sylvain Perifel, and Amir Yehudayoff.
\newblock Separating multilinear branching programs and formulas.
\newblock In Howard~J. Karloff and Toniann Pitassi, editors, {\em Proceedings
  of the 44th Symposium on Theory of Computing Conference, {STOC} 2012, New
  York, NY, USA, May 19 - 22, 2012}, pages 615--624. {ACM}, 2012.

\bibitem[FS13]{ForbesS13}
Michael~A. Forbes and Amir Shpilka.
\newblock Quasipolynomial-time identity testing of non-commutative and
  read-once oblivious algebraic branching programs.
\newblock In {\em 54th Annual {IEEE} Symposium on Foundations of Computer
  Science, {FOCS} 2013, 26-29 October, 2013, Berkeley, CA, {USA}}, pages
  243--252. {IEEE} Computer Society, 2013.

\bibitem[FSS14]{ForbesSS14}
Michael~A. Forbes, Ramprasad Saptharishi, and Amir Shpilka.
\newblock Hitting sets for multilinear read-once algebraic branching programs,
  in any order.
\newblock In David~B. Shmoys, editor, {\em Symposium on Theory of Computing,
  {STOC} 2014, New York, NY, USA, May 31 - June 03, 2014}, pages 867--875.
  {ACM}, 2014.

\bibitem[GG20]{guo2020improved}
Zeyu Guo and Rohit Gurjar.
\newblock Improved explicit hitting-sets for roabps.
\newblock In {\em Approximation, Randomization, and Combinatorial Optimization.
  Algorithms and Techniques (APPROX/RANDOM 2020)}. Schloss
  Dagstuhl-Leibniz-Zentrum f{\"u}r Informatik, 2020.

\bibitem[GKS17]{GurjarKS17}
Rohit Gurjar, Arpita Korwar, and Nitin Saxena.
\newblock Identity testing for constant-width, and any-order, read-once
  oblivious arithmetic branching programs.
\newblock {\em Theory Comput.}, 13(1):1--21, 2017.

\bibitem[GKST17]{GurjarKST17}
Rohit Gurjar, Arpita Korwar, Nitin Saxena, and Thomas Thierauf.
\newblock Deterministic identity testing for sum of read-once oblivious
  arithmetic branching programs.
\newblock {\em Comput. Complex.}, 26(4):835--880, 2017.

\bibitem[GR21]{GR21}
Purnata Ghosal and B.~V.~Raghavendra Rao.
\newblock Limitations of sums of bounded read formulas and abps.
\newblock In Rahul Santhanam and Daniil Musatov, editors, {\em Computer Science
  - Theory and Applications - 16th International Computer Science Symposium in
  Russia, {CSR} 2021, Sochi, Russia, June 28 - July 2, 2021, Proceedings},
  volume 12730 of {\em Lecture Notes in Computer Science}, pages 147--169.
  Springer, 2021.

\bibitem[GV20]{GurjarV20}
Rohit Gurjar and Ben~Lee Volk.
\newblock Pseudorandom bits for oblivious branching programs.
\newblock {\em {ACM} Trans. Comput. Theory}, 12(2):8:1--8:12, 2020.

\bibitem[Jan08]{Jansen08}
Maurice~J. Jansen.
\newblock Lower bounds for syntactically multilinear algebraic branching
  programs.
\newblock In Edward Ochmanski and Jerzy Tyszkiewicz, editors, {\em Mathematical
  Foundations of Computer Science 2008, 33rd International Symposium, {MFCS}
  2008, Torun, Poland, August 25-29, 2008, Proceedings}, volume 5162 of {\em
  Lecture Notes in Computer Science}, pages 407--418. Springer, 2008.

\bibitem[Kay12]{Kayal12}
Neeraj Kayal.
\newblock An exponential lower bound for the sum of powers of bounded degree
  polynomials.
\newblock {\em Electron. Colloquium Comput. Complex.}, {TR12-081}, 2012.

\bibitem[KLSS17]{KayalLSS17}
Neeraj Kayal, Nutan Limaye, Chandan Saha, and Srikanth Srinivasan.
\newblock An exponential lower bound for homogeneous depth four arithmetic
  formulas.
\newblock {\em {SIAM} J. Comput.}, 46(1):307--335, 2017.

\bibitem[KNS20]{KayalNS20}
Neeraj Kayal, Vineet Nair, and Chandan Saha.
\newblock Separation between read-once oblivious algebraic branching programs
  (roabps) and multilinear depth-three circuits.
\newblock {\em {ACM} Trans. Comput. Theory}, 12(1):2:1--2:27, 2020.

\bibitem[KS17]{/KumarS17}
Mrinal Kumar and Shubhangi Saraf.
\newblock On the power of homogeneous depth 4 arithmetic circuits.
\newblock {\em {SIAM} J. Comput.}, 46(1):336--387, 2017.

\bibitem[KS22]{KushS22}
Deepanshu Kush and Shubhangi Saraf.
\newblock Improved low-depth set-multilinear circuit lower bounds.
\newblock In Shachar Lovett, editor, {\em 37th Computational Complexity
  Conference, {CCC} 2022, July 20-23, 2022, Philadelphia, PA, {USA}}, volume
  234 of {\em LIPIcs}, pages 38:1--38:16. Schloss Dagstuhl - Leibniz-Zentrum
  f{\"{u}}r Informatik, 2022.

\bibitem[KS23]{KushS23}
Deepanshu Kush and Shubhangi Saraf.
\newblock Near-optimal set-multilinear formula lower bounds.
\newblock In Amnon Ta{-}Shma, editor, {\em 38th Computational Complexity
  Conference, {CCC} 2023, July 17-20, 2023, Warwick, {UK}}, volume 264 of {\em
  LIPIcs}, pages 15:1--15:33. Schloss Dagstuhl - Leibniz-Zentrum f{\"{u}}r
  Informatik, 2023.

\bibitem[KSS14]{KayalSS14}
Neeraj Kayal, Chandan Saha, and Ramprasad Saptharishi.
\newblock A super-polynomial lower bound for regular arithmetic formulas.
\newblock In David~B. Shmoys, editor, {\em Symposium on Theory of Computing,
  {STOC} 2014, New York, NY, USA, May 31 - June 03, 2014}, pages 146--153.
  {ACM}, 2014.

\bibitem[KST16]{KayalST16}
Neeraj Kayal, Chandan Saha, and S{\'{e}}bastien Tavenas.
\newblock An almost cubic lower bound for depth three arithmetic circuits.
\newblock In Ioannis Chatzigiannakis, Michael Mitzenmacher, Yuval Rabani, and
  Davide Sangiorgi, editors, {\em 43rd International Colloquium on Automata,
  Languages, and Programming, {ICALP} 2016, July 11-15, 2016, Rome, Italy},
  volume~55 of {\em LIPIcs}, pages 33:1--33:15. Schloss Dagstuhl -
  Leibniz-Zentrum f{\"{u}}r Informatik, 2016.

\bibitem[LST21]{LST1}
Nutan Limaye, Srikanth Srinivasan, and S{\'{e}}bastien Tavenas.
\newblock Superpolynomial lower bounds against low-depth algebraic circuits.
\newblock In {\em 62nd {IEEE} Annual Symposium on Foundations of Computer
  Science, {FOCS} 2021, Denver, CO, USA, February 7-10, 2022}, pages 804--814.
  {IEEE}, 2021.

\bibitem[LST22]{LST3}
Nutan Limaye, Srikanth Srinivasan, and S{\'{e}}bastien Tavenas.
\newblock On the partial derivative method applied to lopsided set-multilinear
  polynomials.
\newblock In Shachar Lovett, editor, {\em 37th Computational Complexity
  Conference, {CCC} 2022, July 20-23, 2022, Philadelphia, PA, {USA}}, volume
  234 of {\em LIPIcs}, pages 32:1--32:23. Schloss Dagstuhl - Leibniz-Zentrum
  f{\"{u}}r Informatik, 2022.

\bibitem[MU05]{mitzenmacher_upfal_2005}
Michael Mitzenmacher and Eli Upfal.
\newblock {\em Probability and Computing: Randomized Algorithms and
  Probabilistic Analysis}.
\newblock Cambridge University Press, 2005.

\bibitem[Nis91]{Nisan91}
Noam Nisan.
\newblock Lower bounds for non-commutative computation (extended abstract).
\newblock In Cris Koutsougeras and Jeffrey~Scott Vitter, editors, {\em
  Proceedings of the 23rd Annual {ACM} Symposium on Theory of Computing, May
  5-8, 1991, New Orleans, Louisiana, {USA}}, pages 410--418. {ACM}, 1991.

\bibitem[NW97]{NisanW97}
Noam Nisan and Avi Wigderson.
\newblock Lower bounds on arithmetic circuits via partial derivatives.
\newblock {\em Comput. Complex.}, 6(3):217--234, 1997.

\bibitem[Raz06]{Raz06}
Ran Raz.
\newblock Separation of multilinear circuit and formula size.
\newblock {\em Theory Comput.}, 2(6):121--135, 2006.

\bibitem[Raz13]{Raz-Tensor}
Ran Raz.
\newblock Tensor-rank and lower bounds for arithmetic formulas.
\newblock {\em J. {ACM}}, 60(6):40:1--40:15, 2013.

\bibitem[RR20]{RR20}
C.~Ramya and B.~V.~Raghavendra Rao.
\newblock Lower bounds for special cases of syntactic multilinear abps.
\newblock {\em Theor. Comput. Sci.}, 809:1--20, 2020.

\bibitem[RS05]{RazS05}
Ran Raz and Amir Shpilka.
\newblock Deterministic polynomial identity testing in non-commutative models.
\newblock {\em Comput. Complex.}, 14(1):1--19, 2005.

\bibitem[RY08]{RazY08}
Ran Raz and Amir Yehudayoff.
\newblock Balancing syntactically multilinear arithmetic circuits.
\newblock {\em Comput. Complex.}, 17(4):515--535, 2008.

\bibitem[RY09]{RazY09}
Ran Raz and Amir Yehudayoff.
\newblock Lower bounds and separations for constant depth multilinear circuits.
\newblock {\em Comput. Complex.}, 18(2):171--207, 2009.

\bibitem[Sap15]{Saptarishi-survey}
Ramprasad Saptharishi.
\newblock A survey of lower bounds in arithmetic circuit complexity.
\newblock {\em Github Survey}, 2015.

\bibitem[SY10]{ShpilkaY10}
Amir Shpilka and Amir Yehudayoff.
\newblock Arithmetic circuits: {A} survey of recent results and open questions.
\newblock {\em Found. Trends Theor. Comput. Sci.}, 5(3-4):207--388, 2010.

\bibitem[TLS22]{LST2}
S{\'{e}}bastien Tavenas, Nutan Limaye, and Srikanth Srinivasan.
\newblock Set-multilinear and non-commutative formula lower bounds for iterated
  matrix multiplication.
\newblock In Stefano Leonardi and Anupam Gupta, editors, {\em {STOC} '22: 54th
  Annual {ACM} {SIGACT} Symposium on Theory of Computing, Rome, Italy, June 20
  - 24, 2022}, pages 416--425. {ACM}, 2022.

\bibitem[Val79]{Val79}
L.~G. Valiant.
\newblock Completeness classes in algebra.
\newblock In {\em Proceedings of the Eleventh Annual ACM Symposium on Theory of
  Computing}, STOC '79, page 249–261, New York, NY, USA, 1979. Association
  for Computing Machinery.

\end{thebibliography}
\end{document}
